\chapter{Tutorial: ubiquitin from download to analyzed simulation}
\label{ch:tutorial}

This chapter walks you through one complete project, from downloading a
protein to looking at the results of a short simulation. You will use
ubiquitin (\pdb{1UBQ}), a small, well-behaved protein of 76 amino acids. Every
answer shown here is the answer that was typed during a real run of ProPrep
1.23.0, and every screen is the screen it printed.

Plan on about 30~minutes of your own time. The computer work itself is short:
building the system takes a few seconds and the shortened simulation about one
minute on a laptop.

\begin{glance}
\glanceitem{What you do}{Download \pdb{1UBQ}, surround it with water and salt,
  run a short simulation, and measure how the protein moved.}
\glanceitem{Tools used}{\key{1a} Structure Loader, \key{5a} Topology Generator,
  \key{5b} Molecular Dynamics Manager.}
\glanceitem{Before you start}{ProPrep installed (Chapter~\ref{ch:installation})
  and an Internet connection.}
\glanceitem{Time}{About 30~minutes of reading and typing; about one minute of
  simulation.}
\glanceitem{You get}{\file{1UBQ.prmtop}, \file{1UBQ.rst7}, a folder of
  simulation results, and a table of measurements (\file{.csv}).}
\end{glance}

\section{The plan}

A \term{molecular dynamics} simulation (MD for short) lets the computer move
every atom of a protein, step by tiny step, according to physical rules. You
will go through four stages:

\begin{enumerate}
  \item \textbf{Get the structure.} Download the atom positions of ubiquitin
        from the Protein Data Bank.
  \item \textbf{Build the system.} Choose the physical rules (the force
        field), put the protein in a box of water, and add salt.
  \item \textbf{Simulate.} Relax the system, warm it up, then let it move
        freely while ProPrep records snapshots.
  \item \textbf{Analyze.} Measure how much the protein moved and changed shape.
\end{enumerate}

Some optional steps, such as removing crystal waters or checking protonation
states, are left out here because ubiquitin does not need them. They are
covered in Chapters~\ref{ch:cleaning} and~\ref{ch:protonation}.

\section{Before you begin}\label{sec:tutorial-before}

ProPrep uses the simulation programs of \term{AmberTools} (a free package of
simulation programs that is installed together with ProPrep). If another copy
of Amber is also installed on your computer, ProPrep may quietly use that copy
instead of its own.

\begin{watchout}
In version 1.23.0, ProPrep keeps whatever Amber installation the setting
\texttt{AMBERHOME} already points to. On our test computer it pointed to a
different Amber installation, and every program came from there without any
message. If you have ever installed Amber separately, open a new Terminal
window, start the ProPrep environment as described in
Chapter~\ref{ch:installation}, and check:
\begin{terminal}
echo $AMBERHOME
\end{terminal}
If this prints a folder that does not belong to ProPrep, clear it before
starting ProPrep:
\begin{terminal}
unset AMBERHOME
\end{terminal}
\end{watchout}

Make an empty folder for your projects, move into it, and start ProPrep in Full
Menu mode. In the Terminal:

\begin{terminal}
mkdir ~/proprep-tutorial
cd ~/proprep-tutorial
proprep --full-menu
\end{terminal}

If you prefer the browser version, start \texttt{proprep-web} instead
(Chapter~\ref{ch:starting}). The screens are the same. If ProPrep asks
\prompt{Select menu mode [1/2]}, type \key{2} (Full Menu), and answer \key{n}
when it asks whether to save this as your default.

\section{Step 1: Name the project}

ProPrep first asks for a project name. It makes a folder with that name and
saves every file of the project there.

\begin{screen}
═══ Project Setup ═══

Enter a name for your project. All structures and outputs
will be saved in a directory with this name.
Press Enter to use the current directory.

Project name (.): §ans«ubq»
✓ Project directory: /Users/you/proprep-tutorial/ubq

Session recording to: ./proprep_session_20261121_122842.json
\end{screen}

Type \typed{ubq} and press \key{Enter}. ProPrep also starts a
\term{session file}, a record of every answer you give. It is called
\file{proprep_session_<date>_<time>.json}. You can ignore it for now; Chapter~\ref{ch:sessions}
explains how to reuse it.

The main menu appears next. At this point most tools show a grey circle
(\texttt{○}) and a yellow warning such as \texttt{⚠}\,\prompt{Needs a loaded
structure}.
They become available once a structure is loaded.

\section{Step 2: Download ubiquitin}

\begin{steps}
  \item At \prompt{Enter your choice (1a)}, type \key{1a} (Structure Loader).
  \item Type \key{1} (Load protein structure).
  \item At \prompt{Select source}, type \key{1} (Query Protein Data Bank (RCSB)).
  \item At \prompt{Select method}, type \key{1} (Enter PDB ID directly).
  \item Type \typed{1UBQ}.
  \item At \prompt{Select structure type}, type \key{1} (Asymmetric unit).
\end{steps}

The \term{asymmetric unit} is the structure exactly as it was deposited.
\pdb{1UBQ} contains a single ubiquitin chain, so this is what you want.

\begin{screen}
Enter PDB ID (e.g., 1ABC, 7XYZ): §ans«1UBQ»

Structure type:
  1. Asymmetric unit (default - as deposited)
  2. Biological assembly (functional oligomeric form)
Select structure type (1): §ans«1»

Downloading 1UBQ...

✓ Downloaded 1UBQ

PDB ID: 1UBQ
Classification: CHROMOSOMAL PROTEIN
Title: STRUCTURE OF UBIQUITIN REFINED AT 1.8 ANGSTROMS RESOLUTION
Experimental Method: X-RAY DIFFRACTION
Resolution: 1.8 Å
...
═══ Load Structure - Select Source ═══
...
6. Done (return to the Structure Loader menu)

Select source (1): §ans«6»
\end{screen}

\begin{steps}[resume]
  \item Type \key{6} (Done), then \key{m} to go back to the main menu.
\end{steps}

\textbf{New file:} \file{ubq/1UBQ.pdb}, the downloaded structure. It still
contains the 58 water molecules seen in the crystal. They are kept in this
tutorial; tLEaP treats them like the other water it adds.

\section{Step 3: Choose the force field, water and salt}\label{sec:tutorial-leapin}

Next you describe the simulation system. A \term{force field} is the set of
rules and numbers that says how atoms attract and repel each other. ProPrep
writes your choices into an input file for \term{tLEaP}, the AmberTools
program that builds the system.

\begin{steps}
  \item In the main menu, type \key{5a} (Topology Generator).
  \item In the Topology Generator menu, type \key{2} (Generate tLEaP input for
        single state).
\end{steps}

\begin{screen}
══ Topology Generator Menu ══
  1. ✓ Edit redox site bond definitions
  2. ✓ Generate tLEaP input for single state
  3. ○ Generate tLEaP inputs for all redox microstates
       ⚠ Need batch microstate metadata from Redox Site Prep
  4. ○ Generate prmtop/rst7 files from tLEaP input files
       ⚠ Need to generate tLEaP input first
...
Enter your choice (1): §ans«2»
\end{screen}

ProPrep shows a short panel that explains what tLEaP does, then asks about
force fields one at a time.

\subsection{The protein force field}

\begin{screen}
━━━ PROTEIN FORCEFIELD ━━━
  1     None                          No protein in system
  2     ff19SB (Recommended)          Latest AMBER protein FF - use with OPC water
  3     ff14SB                        Widely validated, good for most applications
...
Select protein forcefield (2): §ans«»
\end{screen}

Press \key{Enter} to accept \textbf{ff19SB}, Amber's current force field for
proteins. At \prompt{Select modified amino acids ()} press \key{Enter} as well:
ubiquitin has no unusual amino acids.

\subsection{Force fields you do not need}

ProPrep then asks for DNA, RNA, carbohydrate, lipid and small-molecule force
fields. Ubiquitin contains none of these, so type \key{1} (None) at each of the
five questions:

\begin{screen}
━━━ DNA FORCEFIELD ━━━
NOTE: For protein+DNA, only OL24 or OL21 are compatible with ff19SB!
  1    None                  No DNA in system
  2    OL24 (Recommended)    Latest with sugar pucker refinements (2024)
...
Select dna forcefield (2): §ans«1»
\end{screen}

\begin{tip}
The default for each of these questions is a real force field, not
\prompt{None}. Pressing \key{Enter} does no harm for a protein (in our recorded
run we pressed \key{Enter} for DNA, RNA, carbohydrates and small molecules), but
the extra files are then loaded for nothing and the input file is longer.
\end{tip}

\subsection{The water model}

The \term{water model} is the recipe used for each water molecule. In the
recorded run we chose \textbf{TIP3P} (\key{6}), the traditional three-point
model, which is fast and widely used.

\begin{screen}
━━━ WATER MODEL ━━━
Must be compatible with protein forcefield!
  1     None                                  Implicit solvent (no water in system)
  2     OPC (Recommended — matches ff19SB)    4-point, optimal for ff19SB
  3     OPC3                                  3-point OPC (faster, good accuracy)
...
  6     TIP3P                                 3-point, traditional (not recommended)
...
Select water model (2): §ans«6»
\end{screen}

ProPrep recommends OPC with ff19SB, and pressing \key{Enter} gives you OPC.
OPC is a four-point model, so each water has one more particle and the
simulation costs a little more. Chapter~\ref{ch:topology} compares the
choices. ProPrep does not warn you when you pick a water model other than the
recommended one.

\subsection{Ion parameters}

\begin{screen}
━━━ DIVALENT+ ION PARAMETERS ━━━
For Mg2+, Ca2+, Zn2+, Fe2+/3+, etc. Monovalent ions auto-loaded with water.
...
  6     12-6-4 TIP3P (Recommended — matches TIP3P)    Best for divalent+ with TIP3P
...
Select divalent+ ion parameters (6): §ans«»
\end{screen}

Press \key{Enter}. The recommended choice always follows the water model you
picked. ProPrep then prints a summary of the selected force fields and a note
that \prompt{No transformer information found}. That note is normal for a
protein without special groups.

\subsection{Water box and salt}

Now ProPrep asks how to fill the box around the protein. Press \key{Enter} at
each question to accept the recommended answer:

\begin{screen}
Solvent Model
  1. Explicit solvent (solvation box with water molecules)
  2. Implicit solvent (GB/PB continuum model, no water box)
Select solvent model [1/2] (1): §ans«»

Box Shape
  1. Truncated octahedron (Recommended) - ~29% fewer waters than rectangular
  2. Rectangular box
Select box shape [1/2] (1): §ans«»

Buffer Distance
Distance from solute to box edge. 8-10 A typical, 12-15 A for large
conformational changes.
Buffer distance (Angstroms) (10.0): §ans«»
\end{screen}

\begin{itemize}
  \item \textbf{Explicit solvent} means real water molecules surround the
        protein, which is the most realistic choice.
  \item A \term{truncated octahedron} is a box with the corners cut off. It
        fits a roughly round protein with less water, so the simulation runs
        faster.
  \item The \term{buffer} of 10\,Å is the smallest distance between the
        protein and the edge of the box. It keeps the protein from feeling its
        own copy across the box edge.
\end{itemize}

Then choose the salt. Press \key{Enter} three times for Na\textsuperscript{+},
Cl\textsuperscript{--} and 150\,mM, the salt level of body fluids. Press
\key{Enter} once more at \prompt{Add another salt? [y/n] (n)}.

\begin{screen}
NaCl concentration
  1. 150 mM NaCl (Recommended) - physiological ionic strength
  2. 100 mM NaCl
  3. 50 mM NaCl
  4. Neutralize only (0 mM)
  5. Custom concentration
Select concentration [1/2/3/4/5] (1): §ans«»

Configured salts:
  1. 150 mM NaCl
Add another salt? [y/n] (n): §ans«»
\end{screen}

(In the recorded run we typed \key{5} and then \typed{150}, which gives the
same 150\,mM.)

ProPrep now runs tLEaP briefly to count the water molecules and the charge of
the protein, and works out how many ions to add:

\begin{screen}
Solvation: truncated octahedron, 10.0 A buffer, 150 mM NaCl

Running tLEaP info pass for accurate ion counts...
  → Net charge: +0.0
  → Water count: 4792
  ✓ Got required info, terminating tLEaP early
...
│  Result: total charge = system (+0) + ions (+0) = 0  ->  13 Na+ + 13 Cl-
...
\end{screen}

Ubiquitin has no net charge, so ProPrep adds 13 sodium and 13 chloride ions,
which gives 150\,mM in 4,792 waters.

\subsection{The input file}

ProPrep shows the tLEaP input it is about to write (shortened here):

\begin{screen}
Current tLEaP Template:
────────────────────────────────────────────────────────────
   1 # ProPrep-generated tLEaP Input File
   2
   3 # === STANDARD FORCEFIELDS ===
   4 source leaprc.protein.ff19SB
...
   9 source leaprc.water.tip3p
  10 loadamberparams frcmod.ions1lm_1264_tip3p
...
  20 mol = loadpdb /Users/you/proprep-tutorial/ubq/1UBQ.pdb
...
  26 solvateoct mol TIP3PBOX 10.00
...
  30 addionsrand mol Na+ 13 Cl- 13
...
  36 saveamberparm mol 1UBQ.prmtop 1UBQ.rst7
...
Edit the tLEaP template? [y/n] (y): §ans«n»
\end{screen}

Type \key{n}. You do not need to change anything.

\subsection{The bond question: answer \texttt{y}}

ProPrep now looks for special bonds, such as bonds to a metal or between two
cysteines. Ubiquitin has none, and ProPrep asks whether to go on without them.

\begin{screen}
STEP 3: tLEaP Input File Generation

No bond definitions found. Run comprehensive RedoxSite detection first.
Do you want to continue without bond definitions? [y/n] (n): §ans«y»
Continuing with empty bond definitions.
Enter output filename for tLEaP input (1UBQ_tleap.in): §ans«»

Successfully wrote tLEaP input file to: 1UBQ_tleap.in
tLEaP template configuration complete!
View the tLEaP input file? [y/n] (y): §ans«»
\end{screen}

\begin{watchout}
Type \key{y} here. In version 1.23.0 the default answer is \key{n}, and
\key{n} stops the whole step: ProPrep prints \prompt{Please run RedoxSite
detection and try again.} and returns to the Topology Generator menu without
writing a file. Your force-field and salt choices are remembered, so if this
happens, type \key{2} again, press \key{Enter} at \prompt{Reconfigure all
settings\ldots? (n)}, answer \key{n} to the edit question and \key{y} here.
\end{watchout}

Press \key{Enter} at the file name and again to view the file.

\textbf{New file:} \file{1UBQ_tleap.in}, the instructions for tLEaP. The
Topology Generator menu now marks option~2 with a dot (\texttt{●}, done) and option~4 is
available.

\section{Step 4: Build the topology}\label{sec:tutorial-build}

Now tLEaP builds the system. It writes two files. The \term{topology}
(\file{.prmtop}) lists every atom, its charge, and every bond and angle with
its force-field numbers. The \term{coordinates} (\file{.rst7}) give the
position of every atom.

\begin{steps}
  \item In the Topology Generator menu, type \key{4} (Generate prmtop/rst7
        files from tLEaP input files).
  \item At \prompt{Run 1UBQ\_tleap.in? [y/n] (y)}, press \key{Enter}.
\end{steps}

\begin{screen}
Running tLEaP for 1UBQ_tleap.in...
TER record issues found in 1UBQ.pdb:
  • Missing TER record between HETATM HOH 77 and HOH 78 in chain A at line 912
...
✓ Inserted 57 TER record(s) in 1UBQ.pdb
Running: tleap -s -f 1UBQ_tleap.in
✓ tLEaP completed successfully

tLEaP Summary: Exiting LEaP: Errors = 0; Warnings = 13; Notes = 1.

Warnings (13):
Warning 1:
  .../bin/teLeap: Warning!
  Close contact of 1.306 angstroms between nonbonded atoms H and H1
...
\end{screen}

Three things happen here:

\begin{itemize}
  \item ProPrep adds the \term{TER records} (end-of-molecule markers) that the
        crystal waters were missing. It changes \file{1UBQ.pdb} itself and
        keeps the original as \file{1UBQ.pdb.ter_backup}.
  \item tLEaP finishes with \textbf{Errors = 0}. That is the number to check.
  \item The 13 warnings all say \prompt{Close contact} between two hydrogen
        atoms. tLEaP has just added the hydrogens, and such small clashes are
        normal; the minimization step in Step~6 removes them.
\end{itemize}

ProPrep then runs \term{ParmEd}, an Amber tool that checks the topology and
adds a small correction for the ions (the 12-6-4 terms you chose with the ion
parameters):

\begin{screen}
Running ParmEd validation on 1UBQ.prmtop...
  ✓ Initial validity check passed
  ✓ Molecule definitions are correct
  Applying 12-6-4 LJ potential (water model: TIP3P)...
  Detected free ions: Cl-, Na+
  ✓ Added LENNARD_JONES_CCOEF (C4 terms) for 12-6-4 potential
  ✓ Saved corrected topology and coordinates
  ✓ Wrote 1UBQ_parmed.pdb (same order as the topology)
ParmEd validation completed successfully

Results: 1/1 files processed successfully
...
│ 1UBQ_tleap.in │ 1UBQ.prmtop     │ ✓ Found       │ 1UBQ.rst7     │ ✓ Found     │
...
All expected topology files were generated successfully!
\end{screen}

The whole step took a few seconds. Your system now has 15,729 atoms: the 76
amino acids of ubiquitin, about 4,800 water molecules and 26 ions.

\textbf{New files} in \file{ubq/}:

\begin{tabularx}{\linewidth}{@{}lL@{}}
\toprule
\textbf{File} & \textbf{What it is} \\
\midrule
\file{1UBQ.prmtop} & The topology (about 2.9\,MB). \\
\file{1UBQ.rst7} & The starting coordinates (about 570\,kB). \\
\file{1UBQ_parmed.pdb} & The whole system as a PDB file, for viewing. \\
\file{1UBQ_parmed_replay.py} & A script that repeats the ParmEd step by hand. \\
\file{leap.log} & tLEaP's full report. \\
\file{1UBQ.pdb.ter_backup} & Your PDB file before the TER records were added. \\
\bottomrule
\end{tabularx}

\begin{tip}
A good topology file is never empty. If something went wrong, ProPrep can
leave a \file{.prmtop} of 0 bytes and still show a tick. Check the size in your
file browser, or in the Terminal with \texttt{ls -l ubq/*.prmtop}.
Chapter~\ref{ch:topology} explains this and other traps.
\end{tip}

Type \key{m} to return to the main menu. \key{5b} now shows a tick.

\section{Step 5: Set up the simulation}\label{sec:tutorial-setup}

The Molecular Dynamics Manager guides you through six short steps, shown in a
bar at the top of each screen:
\prompt{Files → Protocol → Restraints → Queue → Hardware → Review}.

\begin{steps}
  \item In the main menu type \key{5b}, then \key{1} (Setup and configure
        simulations).
\end{steps}

\subsection{Files}

ProPrep finds your topology and coordinates on its own:

\begin{screen}
╭─────────────────────────────────────────────────────────────╮
│ Step 1 of 6: Structure Files (Topology + Coordinates)       │
│ [Files] → Protocol → Restraints → Queue → Hardware → Review │
╰─────────────────────────────────────────────────────────────╯
  Structure Pairs: 1 selected
    1. 1UBQ.prmtop + 1UBQ.rst7
...
Navigation:
   n. Next step (Protocol Selection)
   h. Help
   x. Exit setup

Enter choice: §ans«n»
\end{screen}

Type \key{n} to go on. (This prompt has no default, so \key{Enter} alone gives
an error message; type the letter.)

\subsection{Protocol: start from the Workshop protocol}

A \term{protocol} is the list of simulation stages and their settings. Type
\key{1} (Use a predefined protocol), then \key{4} (Workshop Protocol):

\begin{screen}
Available Protocols:

  1. Membrane Protein Minimization Protocol (2 steps)
  2. Basic Equilibration Protocol (6 steps)
  3. Staged Equilibration at Physiological Conditions (310 K) (9 steps)
  4. Workshop Protocol (CPU, under an hour) (4 steps)
  5. York Lab Equilibration Tutorial (14 steps)
...
Your selection [1/2/3/4/5/b]: §ans«4»
...
Protocol Steps:
  1. Energy Minimization (minimization)
     1000 cycles with the backbone restrained (10 kcal/mol·Å²).
  2. System Heating (heating)
     0 to 300 K over 15 ps, then 5 ps at 300 K, backbone restrained (1.0 ...
  3. NPT Equilibration (equilibration)
     100 ps at 1 bar to settle the density, weak backbone restraints (0.1 ...
  4. NVT Production MD (production)
     300 ps unrestrained, one frame per ps.

Total simulation time: 420.0 ps
\end{screen}

The four stages mean:
\begin{description}
  \item[Minimization] Moves atoms slightly to remove clashes, without any
        motion in time.
  \item[Heating] Warms the system from 0 to 300\,K (room temperature) while
        the protein backbone is held gently in place.
  \item[NPT equilibration] Lets the box size settle at normal pressure so the
        water has the right density.
  \item[Production] The free simulation whose snapshots you analyze.
\end{description}

ProPrep's own note says this protocol takes about 75~minutes with
\texttt{sander} on a laptop. For this tutorial you shorten it to about a
minute. Type \key{e} (Edit protocol) and press \key{Enter} to accept the
suggested protocol name. The workflow editor opens:

\begin{screen}
Editing Workflow: Protocol_1UBQ_1UBQ_20261121_1243
Based on Workshop Protocol (CPU, under an hour)

Workflow Steps:
  1. Energy Minimization (minimization)
  2. System Heating (heating)
  3. NPT Equilibration (equilibration)
  4. NVT Production MD (production)

Commands:
  edit <step>              Edit MDIN parameters (e.g., 'edit 2')
...
  remove <step>            Remove step (e.g., 'remove 5')
...
  save                     Save and finish
...
Your command: §ans«edit 1»
\end{screen}

The settings of each stage are stored in an Amber input file called an
\term{mdin} file. You change them in a simple text editor.

\begin{steps}
  \item \textbf{Shorten the minimization.} Type \typed{edit 1}, then
        \typed{params}. The step's mdin file opens in the text editor
        (\texttt{nano}). Change the two lines
        \texttt{ncyc=500} and \texttt{maxcyc=1000} to
        \begin{screen}
  ncyc=100,
  maxcyc=300,
\end{screen}
        Then press \key{Ctrl-X}, \key{y} and \key{Enter} to save and close
        the editor. ProPrep says \prompt{MDIN parameters updated}. Type
        \typed{back}.
  \item \textbf{Shorten the heating.} Type \typed{edit 2}, then
        \typed{params}. Make the \texttt{\&cntrl} and \texttt{\&wt} sections
        read as below (1,000 steps of 2\,fs is 2\,ps; the temperature rises
        over the first 750 steps and then holds), save, and type
        \typed{back}.
        \begin{screen}
&cntrl
  ntt=3, tempi=0.0, temp0=300.0, gamma_ln=5.0, ig=-1,
  ntc=2, ntf=2, ntb=1, cut=10.0,
  dt=0.002, nstlim=1000,
  ntwr=500, ntwx=50, ntpr=50, ioutfm=1, iwrap=1,
  ntr=1, restraint_wt=1.0, restraintmask='@CA,C,O,N&!:WAT',
  nmropt=1,
/
&wt TYPE='TEMP0', ISTEP1=0, ISTEP2=750, VALUE1=0.0, VALUE2=300.0, /
&wt TYPE='TEMP0', ISTEP1=751, ISTEP2=1000, VALUE1=300.0, VALUE2=300.0, /
&wt TYPE='END' /
\end{screen}
  \item \textbf{Drop the NPT step.} Type \typed{remove 3} and answer \key{y}
        at \prompt{Remove 'NPT Equilibration'? (n)}. The production step
        becomes step~3.
  \item \textbf{Shorten the production.} Type \typed{edit 3}, then
        \typed{params}. Set \texttt{nstlim=500} (500 steps, 1\,ps),
        \texttt{ntwr=250}, \texttt{ntwx=25} and \texttt{ntpr=25} (save a
        snapshot and an energy line every 25 steps). Save and type
        \typed{back}.
  \item Type \typed{save}.
\end{steps}

\begin{watchout}
ProPrep does not check what you type in the editor. During our run a
careless edit deleted the \texttt{imin=1} line of the minimization step and
ProPrep accepted the file without a word. After each edit, read the
\prompt{MDIN Parameters} list that ProPrep prints and check that the lines you
did not mean to change are still there.
\end{watchout}

\begin{tip}
Dropping the NPT step and running only 1\,ps of production is fine for
learning the program. It is far too short for research. For real work keep
all four stages and use the times the protocol suggests
(Chapter~\ref{ch:runningmd}).
\end{tip}

After \typed{save}, ProPrep confirms:

\begin{screen}
Step 2 Complete - 1 protocol(s) configured
...
✓ Added 3 simulations to queue
\end{screen}

\subsection{Restraints and queue}

The next two screens can be left as they are.

\begin{steps}
  \item At \prompt{Step 3 of 6: Restraint Integration}, type \key{s} (Continue
        with current restraints). The minimization and heating steps keep the
        protein backbone gently restrained, as the protocol intends.
  \item At \prompt{Step 4 of 6: Queue Management}, type \key{k} (Keep current
        order and continue).
\end{steps}

\subsection{Hardware: choose sander}\label{sec:tutorial-engine}

The \term{MD engine} is the program that actually runs the simulation.
AmberTools includes \texttt{sander}, which runs on one processor, and
\texttt{sander.MPI}, which can share the work between several processors.
The faster engines \texttt{pmemd} and \texttt{pmemd.cuda} belong to the
separate Amber package, which you obtain and install on your own; they are not
installed with ProPrep.

\begin{screen}
╭─────────────────────────────────────────────────────────────────────╮
│ Step 5 of 6: Hardware Config                                        │
│ Files ✓ → Protocol ✓ → Restraints ✓ → Queue ✓ → [Hardware] → Review │
╰─────────────────────────────────────────────────────────────────────╯
...
  Available: 11 CPUs, 0 GPUs

  Simulations: 3 configured
    3x pmemd

What would you like to do?
   1. Apply recommended assignments (CPU for min/NPT, GPU for rest)
   2. Configure individual simulation engines
   3. Set all simulations to same engine
...
Enter choice (1): §ans«3»

Available engines:
  1. sander
  2. sander.MPI
  3. pmemd
  4. pmemd.MPI
  5. pmemd.cuda
Select engine for all simulations (cancel): §ans«1»
✓ 3 simulations configured with sander
\end{screen}

\begin{steps}
  \item Type \key{3} (Set all simulations to same engine), then \key{1}
        (sander).
\end{steps}

\begin{watchout}
Every step starts out set to \texttt{pmemd}, and the default answer
\key{1} (recommended assignments) picks \texttt{pmemd.cuda} for a graphics
card, even when ProPrep has just reported \prompt{0 GPUs}. Neither program is
installed with ProPrep, and the run then fails with \prompt{No such file or
directory: 'pmemd'}. Always set the engine to \texttt{sander} or
\texttt{sander.MPI} unless you have installed the full Amber package.
\end{watchout}

To make the first two steps faster, we then gave them four processors each.
This is optional:

\begin{steps}[resume]
  \item Type \key{2} to set the engine for each step separately.
  \item For the minimization, type \key{2} (sander.MPI). When ProPrep asks
        for the number of MPI processes, type \typed{4}.
  \item Do the same for the heating.
  \item For the production, press \key{Enter} to keep \texttt{sander}.
\end{steps}

\begin{screen}
Simulation 1: 1UBQ_1UBQ_Energy Minimization
Template: builtin/workshop_cpu/00_minimization.mdin

Current engine: sander

Available engines:
  1. sander (current)
  2. sander.MPI
  3. pmemd (not installed)
  4. pmemd.MPI (not installed)
  5. pmemd.cuda (not installed)
Select engine (skip): §ans«2»
✓ Updated to sander.MPI

MPI Configuration:
Number of MPI processes (1-11) (5): §ans«4»
\end{screen}

This list, unlike the one in option~\key{3}, marks the programs that are not
installed. Type \key{n} to go to the last step.

\subsection{Review and run}

\begin{screen}
                      Simulation Configuration Summary
┏━━━┳━━━━━━━━━━━┳━━━━━━━━━━━━━━━━━━━━━┳━━━━━━━━━━┳━━━━━━━━━━━━━┳━━━━━━━━━━━━┓
┃ # ┃ System    ┃ Simulation          ┃ Ensemble ┃ Sim Time    ┃ Engine     ┃
┡━━━╇━━━━━━━━━━━╇━━━━━━━━━━━━━━━━━━━━━╇━━━━━━━━━━╇━━━━━━━━━━━━━╇━━━━━━━━━━━━┩
│ 1 │ 1UBQ_1UBQ │ Energy Minimization │ —        │ 1000 cycles │ sander.MPI │
│ 2 │           │ System Heating      │ NVT      │ 20 ps       │ sander.MPI │
│ 3 │           │ NVT Production MD   │ NVT      │ 300 ps      │ sander     │
└───┴───────────┴─────────────────────┴──────────┴─────────────┴────────────┘

What would you like to do?
   1. Save setup to workspace (execute later with 'run' command)
   2. Write MDIN only (no SLURM, no launch)
   3. Write files and run locally now
...
Enter choice (1): §ans«3»
\end{screen}

\begin{watchout}
The \prompt{Sim Time} column still shows the protocol's original lengths
(1000 cycles, 20\,ps, 300\,ps). In version 1.23.0 this table does not show
your edits, but the files that are written do contain them.
\end{watchout}

\begin{steps}
  \item Type \key{3} (Write files and run locally now). ProPrep says
        \prompt{Configuration saved, switching to execution\ldots} and opens
        the Execute Simulations menu.
\end{steps}

(In the recorded run we typed \key{2} here, which only writes the files, and
then started the run from \key{5b} \textrightarrow{} \key{2}; the result is
the same.)

\section{Step 6: Run and watch}\label{sec:tutorial-run}

\begin{screen}
===== Execute Simulations =====

Execute Options:
  1. Execute simulation queue (3 simulations)
  2. Manage queue (reorder, hold, remove)

Other Options:
  3. Monitor running simulation
  4. Analyze completed simulations
...
Select option (b): §ans«1»
...
Proceed with execution of 3 active simulations? (y): §ans«y»

Created run directory: .../ubq/simulations/batch_20261121_125742
...
Executing simulation 1/3: 1UBQ_1UBQ_Energy Minimization
Starting simulation in background:
mpirun -np 4 sander.MPI -O -i simulation.mdin -p 1UBQ.prmtop -c 1UBQ.rst7 ...
✓ Simulation '1UBQ_1UBQ_Energy Minimization' started in background (PID: 4080)
...
Executing simulation 2/3: 1UBQ_1UBQ_System Heating
  Queued (depends on: 1UBQ_1UBQ_Energy Minimization)
Will auto-start when dependency completes
...
\end{screen}

\begin{steps}
  \item Type \key{1} (Execute simulation queue) and \key{y}.
\end{steps}

The simulation runs in the background. ProPrep starts each stage as soon as
the previous one has finished, so you can keep using the menus.

\subsection{Watching the progress}

\begin{steps}[resume]
  \item In the Execute Simulations menu, type \key{3} (Monitor running
        simulation), then \key{1} to pick the running step.
  \item Type \key{1} (Show current status) or \key{7} (Plot temperature).
\end{steps}

\begin{screen}
Currently tracked simulations:
  • 1UBQ_1UBQ_System Heating
    (running, runtime: 0:00:10, batch: batch_20261121_125742)
...
Select option (1): §ans«7»

Temperature (K)
============================================================
Max:   291.68                                                ┤
                                               ●  ●  ●  ●    ┤
                                   ●                         ┤
                       ●  ●                                  ┤   145.84
        ●  ●                                                 ┤
●                                                            ┤
Min:     0.00                                                ┤
────────────────────────────────────────────────────────────┘
     0.0                                                 2.0 (ps)

Latest: 291.68 (K) | Average: 161.47 (K) | Points: 21
\end{screen}

The plot (shortened here) shows the temperature climbing from 0 towards
300\,K during heating, exactly as intended. Type \key{b} to leave the monitor.
When all three steps are done, the Execute Simulations menu shows:

\begin{screen}
✓ Finished: Energy Minimization
✓ Finished: System Heating
✓ Finished: NVT Production MD
\end{screen}

These were the times on our test laptop (a Mac with 11 processor cores) for
the 15,729-atom system:

\begin{tabularx}{\linewidth}{@{}lllL@{}}
\toprule
\textbf{Stage} & \textbf{Engine} & \textbf{Length} & \textbf{Time} \\
\midrule
Minimization & sander.MPI, 4 processes & 300 cycles & 6\,s \\
Heating & sander.MPI, 4 processes & 1,000 steps (2\,ps) & 22\,s \\
Production & sander & 500 steps (1\,ps) & 33\,s \\
\bottomrule
\end{tabularx}

\medskip
The minimization brought the energy down from $-26{,}515$ to
$-60{,}537$\,kcal/mol. Heating ended at 291.7\,K, and the production run
averaged 294.5\,K, close to the 300\,K target.

\textbf{New folder:} \file{ubq/simulations/batch_<date>_<time>/}, with one
subfolder per stage and a \file{shared/} folder holding the topology and
starting coordinates:

\begin{screen}
batch_20261121_125742/
  01_1UBQ_1UBQ_Energy_Minimization/
  02_1UBQ_1UBQ_System_Heating/
  03_1UBQ_1UBQ_NVT_Production_MD/
  shared/
\end{screen}

Each stage folder contains:
\begin{itemize}
  \item \file{simulation.mdin}, the settings for the stage;
  \item \file{run_simulation.sh}, the command that was run;
  \item a \file{.mdout} file, the energy report;
  \item a \file{.rst7} file, the final coordinates;
  \item for heating and production, a \file{.nc} file: the
        \term{trajectory}, the file of snapshots that you analyze next.
\end{itemize}

\section{Step 7: Analyze the result}\label{sec:tutorial-analysis}

\begin{steps}
  \item In the Execute Simulations menu, type \key{4} (Analyze completed
        simulations). From the main menu the same tool is \key{5b}, then
        \key{4}.
  \item Press \key{Enter} to choose the newest simulation, the production run.
\end{steps}

\begin{screen}
Available simulations for analysis:
  1. 03_1UBQ_1UBQ_NVT_Production_MD (completed: 2026-11-21 12:58)
  2. 02_1UBQ_1UBQ_System_Heating (completed: 2026-11-21 12:58)
  3. 01_1UBQ_1UBQ_Energy_Minimization (completed: 2026-11-21 12:57)
Select simulation (1-3) or 'browse' (1): §ans«1»
...
Select Analysis Type:
  1. Energetics analysis (mdout)
  2. Trajectory analysis (structural)
  3. Combined analysis (recommended)
  4. View trajectory in the structure viewer
Select analysis type (3): §ans«2»
\end{screen}

\begin{steps}[resume]
  \item Type \key{2} (Trajectory analysis).
  \item A small file browser opens in the stage folder. Type \typed{find},
        then \key{1} to pick the production trajectory,\\
        \file{1ubq_1ubq_nvt_production_md.nc}.
\end{steps}

ProPrep loads the 20 snapshots and shows the analysis menu:

\begin{screen}
        Trajectory Information
┌────────────┬───────────────────────┐
│ Frames     │ 20                    │
│ Atoms      │ 15,729                │
│ Residues   │ 4,926                 │
│ Time range │ 2.0 - 3.0 ps (0.0 ns) │
│ Box type   │ periodic              │
└────────────┴───────────────────────┘
...
══ STRUCTURAL ANALYSIS ══
   1. RMSD - Root mean square deviation
   2. RMSF - Per-residue fluctuations
...
   6. DSSP - Secondary structure (helix, sheet, coil)
...
   9. Distance - Atom-atom distances
...
  21. Radius of gyration - Structural compactness
...
   0. Exit to MD Manager menu

Select analyses (comma-separated, e.g., 1,2,6): §ans«1»
\end{screen}

\subsection{How far did the protein move? (RMSD)}

The \term{RMSD} (root mean square deviation) is the average distance, in Å,
between the atoms in each snapshot and the same atoms in a reference
snapshot. It tells you how far the protein has drifted from where it started.

\begin{steps}[resume]
  \item Type \key{1} (RMSD).
  \item At \prompt{Select region}, type \key{1} (Backbone atoms).
  \item At \prompt{Select reference frame}, press \key{Enter} (Frame 0).
\end{steps}

\begin{screen}
RMSD Analysis: Backbone atoms
Mask: @C,CA,N,O&!:WAT
Reference: Frame 0

     RMSD Statistics
┌───────┬───────────────┐
│ Mean  │ 0.50 ± 0.18 Å │
│ Range │ 0.00 - 0.65 Å │
│ Final │ 0.65 Å        │
└───────┴───────────────┘
...
Analyze additional regions? (y/n) (n): §ans«n»
...
Export analysis results? (4): §ans«3»
✓ Exported to .../ubq/03_1UBQ_1UBQ_NVT_Production_MD_trajectory_analysis.csv
✓ Exported to .../ubq/03_1UBQ_1UBQ_NVT_Production_MD_trajectory_analysis.json
\end{screen}

The backbone moved 0.65\,Å from its starting shape. For a stable protein over
1\,ps that is what you expect: the fold is intact.

\begin{steps}[resume]
  \item Answer \key{n} to \prompt{Analyze additional regions?}.
  \item At \prompt{Export analysis results?}, type \key{3} (Both formats).
\end{steps}

The \file{.csv} file opens in Excel or any spreadsheet program; the
\file{.json} file is for scripts.

\subsection{More measurements at once}

Back at the analysis menu, type several numbers separated by commas:
\typed{2,9,21,6}. ProPrep runs them one after the other.

\begin{steps}[resume]
  \item \textbf{RMSF} (\key{2}): press \key{Enter} three times (C-alpha atoms,
        default alignment, frame~0).
  \item \textbf{Distance} (\key{9}): type \typed{:1@CA} and \typed{:76@CA}
        (the first and last amino acid), and a description such as
        \typed{N-C terminus}.
  \item \textbf{Radius of gyration} (\key{21}): type \key{4} (Protein residues
        1-76, no hydrogens).
  \item \textbf{DSSP} (\key{6}): type \key{4} as well.
  \item At \prompt{Export analysis results?}, type \key{1} (CSV).
\end{steps}

Our results, in plain words:

\begin{tabularx}{\linewidth}{@{}lrL@{}}
\toprule
\textbf{Measurement} & \textbf{Result} & \textbf{Meaning} \\
\midrule
RMSF & 0.28\,Å mean & How much each amino acid wobbles. The tail
  (residue~76, 0.72\,Å) moved most. \\
Distance & 38.4\,Å & The two ends of the chain stayed about 38\,Å apart. \\
Radius of gyration & 11.9\,Å & The protein stayed equally compact. \\
DSSP & 20\,\% helix, 34\,\% sheet & The helix and sheets of ubiquitin
  were kept. \\
\bottomrule
\end{tabularx}

\medskip
Chapter~\ref{ch:analysis} explains each measurement and how to read it.

\begin{watchout}
Never type a number containing the digit 0, such as \key{10} or \key{20}, in
the analysis list. In version 1.23.0 any answer with a 0 in it is read as
\key{0} (Exit), and ProPrep leaves the analysis menu at once.
\end{watchout}

\begin{steps}[resume]
  \item Type \key{0} to leave the analysis menu, then \key{x} to close
        ProPrep when you are finished.
\end{steps}

\section{What you just did}

\begin{tabularx}{\linewidth}{@{}llL@{}}
\toprule
\textbf{Stage} & \textbf{Keys} & \textbf{Files in \file{ubq/}} \\
\midrule
Named the project & --- & \file{proprep_session_<date>.json} \\
Downloaded the structure & \key{1a} & \file{1UBQ.pdb} \\
Chose force field, water, box, salt & \key{5a} \key{2} & \file{1UBQ_tleap.in} \\
Built the system & \key{5a} \key{4} & \file{1UBQ.prmtop}, \file{1UBQ.rst7},
  \file{leap.log} \\
Set up and ran the simulation & \key{5b} \key{1} & \file{simulations/batch_<date>/} \\
Analyzed the trajectory & \key{5b} \key{4} & \file{03_..._trajectory_analysis.csv} \\
\bottomrule
\end{tabularx}

\medskip
You took a protein from the Protein Data Bank, gave it a force field, put it
in a box of water with physiological salt, relaxed and warmed it, ran it
freely, and measured how it behaved. Every real simulation project follows
the same path; it only adds preparation steps and runs much longer.

\section{If something went differently}

\begin{tabularx}{\linewidth}{@{}lL@{}}
\toprule
\textbf{What you see} & \textbf{What to do} \\
\midrule
Back at the Topology Generator menu, no \file{.in} file &
  You answered \key{n} to the bond question. Run \key{2} again and answer
  \key{y} (Section~\ref{sec:tutorial-leapin}). \\
\prompt{No such file or directory: 'pmemd'} &
  The engine was left at \texttt{pmemd}. Go back to \key{5b} \key{1} and set
  \texttt{sander} at the hardware step (Section~\ref{sec:tutorial-engine}). \\
\prompt{Returning to MD Manager menu\ldots} right after choosing analyses &
  Your list contained a 0. Choose again without 10 or 20. \\
A \file{.prmtop} of 0 bytes &
  tLEaP failed. See Chapter~\ref{ch:topology}, ``If something goes wrong''. \\
\bottomrule
\end{tabularx}

\medskip
More problems and their solutions are in Chapter~\ref{ch:troubleshooting}.

\section{Where to go next}

\begin{itemize}
  \item \textbf{Your own protein:} Chapter~\ref{ch:structure} shows the other
        ways to load a structure, and Chapters~\ref{ch:cleaning}
        and~\ref{ch:protonation} show how to clean it and set protonation
        states.
  \item \textbf{More about the system:} Chapter~\ref{ch:topology} explains
        every choice in the Topology Generator, including disulfide bonds and
        hemes.
  \item \textbf{Longer and better simulations:} Chapter~\ref{ch:runningmd}
        covers protocols, restraints, engines and running on a computing
        cluster.
  \item \textbf{More analyses:} Chapter~\ref{ch:analysis} lists all 23
        measurements.
\end{itemize}
